\documentclass[11pt]{article}
\usepackage[margin=1in]{geometry}
\usepackage{amsmath,amssymb,amsthm,hyperref}
\newtheorem{theorem}{Theorem}
\newtheorem{corollary}[theorem]{Corollary}
\newcommand{\tr}{\operatorname{tr}}
\title{Exact tensor activation thresholds for a five-dimensional reflection}
\author{Joshua Mark}
\date{20 November 2026; research derivation, priority not established}
\begin{document}\maketitle
Put $F_p(U)=\sum_{ij}|u_{ij}|^p$ on $O(n)$, where $p>2$.
A negative tangent Hessian means negative second derivative along
$Ue^{tX}$ for a real skew-symmetric $X$. This note sharpens the finite
tensor activation theorem with exact examples; it does not establish
Hadamard existence or claim priority.

\begin{theorem}\label{main}
Let $A=2J_5/5-I_5$, $s=p-2$, and
\[
 b=2/5,\qquad c=(3/5)^{p-1}+(2/5)^{p-1},\qquad
 R(p)=c/b^s=\tfrac35(\tfrac32)^s+\tfrac25.
\]
For every positive integer $k$, the least eigenvalue of the tangent
Hessian of $F_p$ at $A^{\otimes k}$, with the Frobenius metric on
skew generators, is exactly
\[
 \lambda_k=p\big((p-1)b^{ks}-c^k\big).
\]
Consequently it is a strict local minimum if $R(p)^k<p-1$, has positive
semidefinite singular Hessian if equality holds, and has a negative
tangent direction if $R(p)^k>p-1$. The first negative-Hessian power is
\[
 k_-(p)=\left\lfloor\frac{\log(p-1)}{\log R(p)}\right\rfloor+1.
\]
The equality case does not assert local minimality.
\end{theorem}
\begin{proof}
Write $d=(3/5)^{p-1}$ and $e=(2/5)^{p-1}$. The entrywise gradient
matrix is $W=eJ_5-(d+e)I_5$, so $M=A^TW$ is symmetric and $A$ is
critical. On $\mathbf1^\perp$, $M$ has eigenvalue $c=d+e$, while on
$\mathbf1$ it has eigenvalue $h=-d+4e$. For $p>2$,
$3e<2d$, hence $-c<h<c$. Thus $\|M\|_{\rm op}=c$.
Tensor multiplicativity gives criticality and $M_k=M^{\otimes k}$.
Every entry of $A^{\otimes k}$ has magnitude at least $b^k$.
For any skew $X$, orthogonality and the operator norm bound imply
\[
 T_k(X):=\sum_{ij}|(A^{\otimes k})_{ij}|^s
                 ((A^{\otimes k}X)_{ij})^2\ge b^{ks}\|X\|_F^2,
 \qquad
 S_k(X):=\tr(M_kX^TX)\le c^k\|X\|_F^2.
\]
Here the last bound remains valid if $M_k$ has negative eigenvalues.
The derivative identity $F_p''=p[(p-1)T_k-S_k]$ gives the claimed
lower bound for the least Hessian eigenvalue.

To attain it, take a skew matrix $Y$ supported on the oriented triangle
$0\to1\to2\to0$, with entries $+1$ on these three edges and $-1$
on reversed edges. Take symmetric $Z$ supported on the undirected
cycle $0,1,2,3,0$, with successive edge values $+1,-1,+1,-1$.
Both have zero diagonal and zero row and column sums. Consequently
$AY=-Y$, $AZ=-Z$, and $MY=cY$, $MZ=cZ$.
For $X=Y\otimes Z^{\otimes(k-1)}$, each nonzero entry has distinct
row and column indices in every coordinate. Thus it encounters only
magnitude-$b^k$ entries of $A^{\otimes k}$ in $T_k$, and
\[
 T_k(X)=b^{ks}\|X\|_F^2,\qquad S_k(X)=c^k\|X\|_F^2.
\]
The matrix $X$ is nonzero and skew, proving sharpness. Positive
definite Hessian at a critical point yields a strict local minimum.
\end{proof}

\begin{corollary}[Two copies do not suffice]
For $p=3$, $A$ and $A^{\otimes2}$ are strict local minima, while
$A^{\otimes3}$ has a negative tangent direction. For $p=201/100$,
the first four tensor powers are strict local minima, while the fifth
has a negative tangent direction.
\end{corollary}
\begin{proof}
At $p=3$, $R=13/10$ and $R^2<2<R^3$.
For $p=201/100$, exact integer exponentiation verifies
\[
 (1004062/10^6)^{100}<3/2<(1004063/10^6)^{100}.
\]
Put $a=\frac35(1004062/10^6)+\frac25$ and
$z=\frac35(1004063/10^6)+\frac25$. Exact rational arithmetic gives
$z^4<101/100<a^5$. Hence $R^4<p-1<R^5$.
\end{proof}

The example at $p=3$ refutes universal two-copy negative-Hessian
activation, including when the seed is a strict local minimum.
The first-power dimension here is five, so it is not itself an
admissible Hadamard order. Tensoring with a normalized Hadamard matrix
also preserves this exact Hessian classification: the minimum entry
weight and the operator norm of the critical matrix both acquire
the same positive factor. An attaining generator is the generator
above tensored with the identity on the flat factor. Thus taking
$A\otimes(H_4/2)$ gives the same examples in initial dimension twenty.
For its $k$th self-power the flat factor has order $4^k$; the Hessian
least eigenvalue is multiplied by $(4^k)^{-(p-2)/2}$.

All activation indices in this specific family are at most five.
Indeed convexity of $\log(\frac35e^{s\log(3/2)}+\frac25)$ gives
$\log R\ge\frac35s\log(3/2)$, whereas $\log(1+s)<s$.
Since $3\log(3/2)>1$, it follows that $5\log R>\log(p-1)$.
Moreover the threshold tends to $5/(3\log(3/2))\simeq4.1105$ as
$p\downarrow2$. This proves that five is attained, but does not
settle whether other critical points require arbitrarily many copies.
\end{document}
